\section{Introduction}\label{sec:intro}

Warehouse order picking is usually the single most labour-intensive operation
in a distribution or supply facility.
Travel alone accounts for roughly half of all picking time~\cite{dekoster2007},
so even modest reductions in walking and equipment-travel distances translate
directly into throughput gains and lower operating cost.
Storage location assignment---deciding which stock-keeping unit lives in which
bin---is one of the few design levers that can improve picking efficiency
without capital expenditure~\cite{gu2007,bartholdi2014}.
Policies range from simple ABC frequency-of-use ranking to multi-criteria
optimization approaches, each with different consequences for the physical
flow of pickers and material-handling equipment.

Schneider Electric's Manisa medium-voltage switchgear plant runs a
pallet-rack warehouse that feeds several parallel production lines through
a kitting process.
Kitting here means assembling a complete set of components for one production
order---a kit---before any assembly work begins; the warehouse dispatches not
single parts but coordinated batches of components, one batch per order.
Because every component of a kit must be collected before the kit is released
to the line, preparation time depends directly on where each component is
stored.
The more scattered the bin assignments, the more the picker travels and the
more often a reach truck (RT) must be called to retrieve pallets from upper
levels, adding queuing time on top of travel time.

\subsection{Problem Statement}

The Manisa warehouse currently assigns storage bins through a combined
ABC/FMR (fast-mover/medium-mover/rare-mover) classification tool in SAP
that ranks each material by total consumption, with no distinction for which
production line it feeds.
High-frequency materials end up close to the kitting area in the aggregate,
but materials belonging to the same kit order are often scattered across
different rack rows and levels, forcing the picker to traverse multiple aisles
and repeatedly call for RT assistance at upper-level bins.
That scatter drives up walking distance, increases the number of RT
interventions per order, and creates queuing for the seven reach trucks on
the floor---all of which inflate kitting lead time.
We were therefore asked to evaluate whether alternative storage assignment
policies could reduce these inefficiencies before any physical re-slotting of
the warehouse is undertaken.

\subsection{Motivation}

The warehouse processes roughly 400 kit orders every active working day---the
SAP~ZWM92 dispatch log covering 2026-01-02 to 2026-05-18 records
40\,804 kit orders across 103 active days, giving a mean of about 396
orders per day.
A one-minute reduction in average kitting lead time therefore saves the
equivalent of more than six full working hours of aggregate delay each day,
and that saving recurs every shift.
Beyond lead time, every kit order filled entirely by a floor-level operator
pick avoids an RT intervention, keeping the reach trucks free for replenishment
and inbound put-away---tasks that cannot be deferred the way a kitting trip can.
The economic case for improving storage assignment is clear even without
elaborate modelling: the sheer order volume makes small per-order gains
significant at the plant level.

Physically re-slotting thousands of pallet positions is, however, expensive
and disruptive: bins must be vacated, relabelled, and refilled, often during
production, and a poorly chosen policy could make matters worse.
A validated discrete-event simulation (DES) model lets us test candidate
policies against four months of real demand data without touching a single
pallet~\cite{negahban2014,law2015}.
By driving the model with empirical inter-batch gaps and kit compositions
extracted from the ZWM92 log, and running twenty independent replications per
policy under common random numbers, the simulation produces statistically
comparable estimates of lead time, waiting time, RT utilization, throughput,
and walking distance---evidence that would be impractical to collect from the
live warehouse in any reasonable timeframe.

\subsection{Objectives}

This project pursues four objectives:

\begin{enumerate}
  \item Build a data-driven discrete-event simulation of the Manisa
        warehouse, parametrised from the SAP material master, the ZWM92
        dispatch log, and an on-site F400 video time-motion study.
  \item Validate the simulation model against the historical dispatch log
        using chi-square goodness-of-fit and per-material pick frequency
        comparisons, following the Sargent verification and validation
        framework~\cite{sargent2013}.
  \item Design and simulate six storage assignment policies---the current
        SAP baseline, a capacity heuristic, usage-based ABC, Double
        ABC, a travel-distance optimized assignment, and the proposed
        Line-aware Slotting policy that groups materials by the production
        line they serve---and compare them on the five key performance
        indicators above.
  \item Recommend a placement strategy with statistical evidence, drawing
        on ANOVA, CRN paired $t$-tests, and 95\% confidence intervals
        derived from $N = 20$ replications.
\end{enumerate}
